% !TEX root = ../Thesis.tex

\section{The Premise Removed by Chapter 6}

Chapter 6 leaves two facts standing together. Five independent estimators
recover the
trial's average effect, and none of them finds a per-patient dispersion of usable magnitude. 
Its closing section also gives a
mechanism for the second fact: MASTER DAPT enrolls only patients already flagged at high
bleeding risk, so the bleeding-risk axis is compressed before any model sees a single
patient, and the aggregate ischemic score sums a MI effect and a stroke effect of
opposite sign, so a real cancellation looks identical to a real absence. The same
compression is why bleeding, with 506 events, is predicted worse (AUROC 0.62) than
cardiovascular death, with 81 (AUROC 0.76) in Chapter 5: enrollment flattened precisely
the dimension a bleeding model would need.

\vspace{10pt}

This removes the premise of the strategy Chapter 3 was built for. An enrollment
mechanism aims to select the cohort with the best combined gain in bleeding and ischemic
risk; if the conflict between the two axes is mostly compression and cancellation rather
than a real signal, there is nothing left for the mechanism to select on. The question
this chapter asks is not whether guided enrollment would help MASTER DAPT, but whether
the machinery of Chapter 3, applied with the same criteria, can still be distinguished
from a machinery that selects at random. The estimated CATE is taken at face value,
cohorts are built from it, and what is measured is whether those cohorts differ in
\textbf{observed} bleeding and ischemic events.

\section{Six Policies on the Trade-off Plane}

\paragraph{The mechanism.}
The enrollment mechanism follows the same skeleton as Chapter 3: a seed cohort of 1,000
patients drawn at random, hyperparameters tuned once on that seed by the same
cross-validated AUTOC lower-bound criterion as Chapter 2 and left fixed for the rest of
the run, and periodic refitting on the patients enrolled so far. It differs from every
mechanism introduced there in how a round is decided: rather than sampling a batch and
keeping its single best candidate, every not-yet-enrolled patient is scored with the
current model, the 100 highest-scoring are enrolled, the model is refit on the enlarged
cohort, and the round repeats. Enrollment stops at 4,200 of the 4,579 patients.

\vspace{10pt}

\paragraph{Five policies.}
As previously anticipated, the trade-off plane is built from two indipendent components,
The bleeding CATE and the weighted ischemic CATE.
This multidimensionality allows for several ways to combine them into a single scalar score.
The five policies and the control are explained below.

\begin{itemize}
    \item The \emph{\textbf{win-win score}} combines the two components, prioritising patients in the top-right quadrant, where both bleeding and ischemic outcomes are expected to improve.
    
    \item The \emph{\textbf{trade-off score}} subtracts the ischemic component from the bleeding component, prioritising patients in the bottom-right quadrant, where a reduction in bleeding risk is accompanied by an increase in ischemic risk.
    
    \item The \emph{\textbf{diagonal scores}} uses the geometry of the treatment-effect plane directly. Patients in the win-win quadrant are ranked above those outside it; within each quadrant, ties are resolved according to distance from the origin or proximity to the $+45^\circ$ diagonal.
    
    \item The \emph{\textbf{pure ischemic score}} rank patients according to the ischemic component alone, using the opposite direction as control.
    
    \item The \emph{\textbf{random policy}} enrolls patients without using their estimated CATEs, this policy is our "0" value.
\end{itemize}

\vspace{10pt}

The diagonal-score uses raw coordinates rather than a linear combination
to rank the patients, so it is not a scalar score in the same sense as the others.
Each policy is run over 100 independent seed cohorts, under two batch-sampling
temperatures: one that occasionally accepts a candidate other than the top-scoring one
within a batch, following the exploration rule of Chapter 3's Mechanism 1.

\section{Drift on the Reference Plane}

Whether a policy moves the cohort at all is measured against the frozen reference model
of Chapter 3: the tuned forest of Chapter 6, fitted once on the whole population and
never used to guide selection, scores every policy's final enrolled cohort so that the
yardstick does not move with the thing it measures. Table~\ref{tab:ch8-drift} reports the
result at the 4,200-patient checkpoint, greedy configuration, averaged over 100 seeds.

\begin{table}[h]
\centering
\caption{Mean reference-plane coordinates of the enrolled cohort at $n=4{,}200$, greedy
configuration, averaged over 100 seed cohorts, scored by the frozen tuned forest of
Chapter 6. \texttt{Bleed} and \texttt{isch} are the cohort mean of the bleeding and
weighted-ischemic CATE in IQR-scaled units; win--win share is the fraction of the enrolled
cohort with both raw CATEs positive.}
\label{tab:ch8-drift}
\begin{tabular}{lccc}
\toprule
Policy & Bleed & Isch & Win--win share \\
\midrule
Win--win score      & 4.462 & 0.338 & 79.4\% \\
Trade-off score      & 4.451 & 0.319 & 78.6\% \\
Diagonal geometry     & 4.460 & 0.340 & 79.6\% \\
$-$isch (control)    & 4.417 & 0.303 & 78.0\% \\
$+$isch (control)    & 4.438 & 0.340 & 79.7\% \\
Random               & 4.430 & 0.320 & 78.6\% \\
\midrule
Full cohort           & 4.430 & 0.322 & 78.6\% \\
\bottomrule
\end{tabular}
\end{table}

Every policy moves in the direction its own objective points: the ischemic axis is
highest under $+$isch and the diagonal rule, lowest under $-$isch, and the random arm sits
on the full-cohort baseline in every column, which is the check that validates the
reference model rather than a result in its own right. The movement is nonetheless
small. The ischemic axis spans $0.303$ to $0.340$ against a baseline of $0.322$, a range
of about $\pm6\%$; the bleeding axis and the win--win share move by about one percentage
point at most. At a 4,200-patient checkpoint, 92\% of the cohort is already enrolled by
construction, which leaves little room for any policy to diverge from the population it
is drawn from — a limitation returned to below, not a property of the policies
themselves.

\section{The Mirror Test}

The two single-axis controls exist to answer a question the diagonal rule and the two
composite scores cannot: if the ischemic axis genuinely ranks patients by ischemic
vulnerability, pushing it in one direction should symmetrically move how many real
ischemic-event patients are enrolled early versus deferred. For each of the five
non-random policies, the drift on the ischemic axis ($\Delta_{\text{isch}}$) is regressed
against the shift in the mean enrollment position of the 205 patients who actually had an
ischemic event ($\Delta_{\text{pos}}$), read at an intermediate round of 3,100 enrolled
patients so the checkpoint is not already exhausted. Table~\ref{tab:ch8-mirror} reports
the five points.

\begin{table}[h]
\centering
\caption{Ischemic-axis drift and shift in the mean enrollment position of the 205
real ischemic-event patients, at $n=3{,}100$, tempered configuration, 30-seed average.
A negative $\Delta_{\text{isch}}$ pairs with a positive $\Delta_{\text{pos}}$ if pushing
the axis down genuinely defers ischemic patients.}
\label{tab:ch8-mirror}
\begin{tabular}{lrr}
\toprule
Policy & $\Delta_{\text{isch}}$ & $\Delta_{\text{pos}}$ (patients) \\
\midrule
Win--win score      & $+0.028$ & $+7.5$ \\
Trade-off score      & $-0.014$ & $+63.6$ \\
Diagonal geometry     & $+0.031$ & $-0.9$ \\
$-$isch              & $-0.035$ & $+49.7$ \\
$+$isch              & $+0.036$ & $-33.2$ \\
\bottomrule
\end{tabular}
\end{table}

The tempered configuration gives a Pearson correlation of $r = -0.894$ ($p = 0.016$)
between the two columns, with the sign a mirror test should have: a more negative
ischemic drift pairs with real ischemic patients entering later. The same regression on
the greedy configuration gives $r = -0.114$ ($p = 0.829$), and the fitted intercept
changes sign between the two, $+28$ patients against $-45$. With only five policies
feeding the regression, the sign and the significance of the mirror test depend on which
exploration setting produced the cohort, and neither should be treated as a stable
result.

\vspace{10pt}

What can be said is only the qualitative direction in the tempered configuration, not a
stable magnitude.

\section{Real Outcomes}

The reference plane is built from the same tuned forest as Chapter 6, so agreement with
it is not independent evidence that anything real changed. The observed bleeding and
ischemic events are read directly from the trial's adjudicated outcomes and never pass
through a CATE estimate. For each policy and each of the two temperature
settings, the ischemic event rate of the enrolled cohort is compared with that of the
patients left out, paired over 30 seeds (Table~\ref{tab:ch8-outcomes}).

\begin{table}[h]
\centering
\caption{Paired difference in ischemic event rate, enrolled minus left-out, at
$n=4{,}200$, 30-seed average, standard error of the mean and two-sided $p$-value from a
paired $t$-test. Positive values mean the enrolled cohort has \emph{more} ischemic
events than the patients left out.}
\label{tab:ch8-outcomes}
\begin{tabular}{lrrr|rrr}
\toprule
& \multicolumn{3}{c|}{Tempered} & \multicolumn{3}{c}{Greedy} \\
Policy & diff (pp) & SEM & $p$ & diff (pp) & SEM & $p$ \\
\midrule
$+$isch              & $+0.40$ & 0.25 & 0.116 & $-0.20$ & 0.29 & 0.506 \\
Diagonal geometry     & $+0.19$ & 0.28 & 0.509 & $-0.14$ & 0.30 & 0.644 \\
Win--win score       & $+0.05$ & 0.27 & 0.857 & $-0.34$ & 0.27 & 0.211 \\
Random               & $+0.03$ & 0.24 & 0.903 & $+0.14$ & 0.21 & 0.524 \\
Trade-off score       & $-0.93$ & 0.24 & \textbf{0.0006} & $-1.00$ & 0.22 & \textbf{0.0001} \\
$-$isch              & $-1.14$ & 0.32 & \textbf{0.0012} & $-0.66$ & 0.23 & \textbf{0.0069} \\
\bottomrule
\end{tabular}
\end{table}

The trade-off score and $-$isch are the only two arms that separate from noise, and both
in the same direction, in both temperature settings: their enrolled cohorts carry
\emph{fewer} ischemic events than the patients left out, not more. This is the opposite
of what a mechanism meant to concentrate on the trade-off quadrant would need to show,
and it is consistent with the small, unstable movement of Table~\ref{tab:ch8-drift}
rather than with a real enrichment. The win--win score, the diagonal rule and random
never separate from zero at either temperature.
Figure~\ref{fig:ch8-ischaemic-trajectory} shows the same comparison over the whole
enrollment run rather than at one checkpoint: the enrolled and left-out ischemic rates
track each other closely, within their own standard-error bands, for every policy.

\begin{figure}[h]
\centering
\includegraphics[width=\textwidth]{GuidedEnrollment_ischaemic_patients.png}
\caption{Cumulative count and rate of real ischemic events in the enrolled cohort
(crimson) against the patients still left out (orange), as enrollment proceeds from the
1,000-patient seed to the 4,200-patient checkpoint, greedy configuration, mean of 30
seed cohorts with standard-error bands. The dotted line marks the cohort's own ischemic
event rate. The two bands overlap for every policy at every stage of enrollment.}
\label{fig:ch8-ischaemic-trajectory}
\end{figure}

\vspace{10pt}

A second, model-internal check reaches the same place from inside the estimation itself.
Scoring the withheld arrivals at every round with a doubly-robust AUTOC lower bound, none
of the five non-random policies clears zero in any of the 30 seeds at the 4,000-patient
checkpoint: the best of them, $-$isch and the trade-off score, average $-0.006$, and the
win--win score and the diagonal rule average $-0.007$. The same estimator that Chapter 6
could never separate from its own noise floor does not separate from it here either, even
when it is the one steering enrollment.

\vspace{10pt}

\paragraph{The correlation behind the collapse.}
The win--win score's own construction explains why it moves so little. On the tuned
forest of Chapter 6 it correlates $\rho = +0.91$ with the bleeding CATE alone and only
$\rho = +0.44$ with the weighted ischemic CATE: bleeding's raw effect keeps roughly
double the spread of the ischemic composite even after both are rescaled to a common IQR,
so ranking by the sum mostly re-ranks by bleeding. Of the 1,500 patients it ranks
highest, 58\% fall in the true win--win quadrant, not the near totality the name
suggests. It is the same mechanism Chapter 6 found in the aggregate ischemic score,
looked at from the selection side rather than the estimation side: a linear combination
of two quantities of unequal variance collapses toward the larger one.

\section{The Bandit Mechanisms Against Their Own Ceiling}

Mechanisms 4 and 5 of Chapter 3, UCB1 and Thompson sampling, are also run on this cohort,
over a duel of arriving pairs rather than a whole-pool sweep, with the exploration bonus
active. Two things beyond the enrolled cohort's own view of itself are needed to judge
them honestly: a \emph{referee}, the same frozen tuned forest used above, scoring the
final enrolled cohort with a model that took no part in selecting it; and a
\emph{ceiling}, the identical sequence of arrival pairs scored with the referee instead of
the online model, which is the best either bandit could have done with a perfect forest
and no learning delay.

\begin{table}[h]
\centering
\caption{Referee-scored conflict of the bandit-enrolled cohort against its own ceiling
(same arrival pairs, offline tuned forest) and against a matched random cohort of equal
size ($n = 2{,}789$), plus Fisher's exact test on real bleeding, death, MI and stroke
event rates between the guided cohort and its random match.}
\label{tab:ch8-bandits}
\begin{tabular}{lccc|c}
\toprule
Policy & Online-guided & Ceiling & Fraction of ceiling & Fisher exact, worst $p$ \\
\midrule
UCB1     & $+0.091$ & $+0.303$ & 30\% & 0.349 (death) \\
Thompson & $+0.105$ & $+0.330$ & 32\% & 0.281 (MI) \\
Random   & $\approx 0.000$ & --- & --- & --- \\
\bottomrule
\end{tabular}
\end{table}

Both bandits reach a real, positive gap over random, about a third of the ceiling
reachable on the identical arrivals if the referee itself had been doing the choosing.
The gap is not the bottleneck's whole story: UCB1 is deterministic across all 30 replicate
runs, since its exploration bonus is a fixed function of the current model's uncertainty
with no random component, while Thompson's sampling step gives it run-to-run variation.

\vspace{10pt}

None of it moves real outcomes. Fisher's exact test on the guided cohort against a
matched random cohort of the same size finds no endpoint distinguishable at conventional
significance for either policy; the smallest $p$-value, for Thompson on MI, is $0.281$.
The online model that guides these two mechanisms is deliberately light, refit every 100
enrollments without its own hyperparameter search past the seed, and the gap to the
ceiling is consistent with that: the bottleneck is the forest's own noise under a small,
frequently-refit sample, not a defect specific to either exploration rule.

\section{What This Chapter Establishes}

Every mechanism moves the estimated-CATE plane in the direction its objective points,
and the movement survives a reference model that never sees the selection. It is also
small at the checkpoint examined, at most a few percentage points on any axis, in a
cohort where 92\% of the patients are already enrolled by the time the comparison is
made. On real, adjudicated outcomes, the picture does not improve on a larger checkpoint
or a different exploration setting: the two arms that do separate from a random control
move in the wrong direction for a mechanism meant to concentrate on the trade-off
quadrant, the win--win score and the diagonal rule never separate from noise at either
temperature, the model-internal AUTOC lower bound never clears zero for any of them, and
the bandit mechanisms enrich their own referee-scored score by a real but partial margin
without moving a single real event rate.

\vspace{10pt}

The reference-plane result and the real-outcome result are not two independent
confirmations of the same thing; only the second is. Chapter 6's tuned forest supplies
both the scores that steer enrollment and the yardstick that measures the drift, so
agreement between them shows the machinery is internally consistent, not that it found
something real. What is independent is the observed event counts, which never pass
through an estimated CATE, and they show the same absence of signal Chapter 6 already
established. \emph{The machinery selects well on a signal that is mostly noise.}

\todo{Re-run the whole-pool comparison at an earlier checkpoint (for instance
$n=2{,}500$--$3{,}000$): at 92\% enrollment the policies have too little of the cohort left
to diverge from, which biases the comparison toward the null already reported here.}
\todo{Build a matched-size, non-adaptive randomized-arrival baseline so that precision, not
only cohort composition, can be compared between guided and plain enrollment at equal
enrolled $n$; no such comparison currently exists for this cohort.}
